\section*{CHAPTER 2: SYSTEM ANALYSIS AND DESIGN}
\addcontentsline{toc}{section}{\numberline{}CHAPTER 2: SYSTEM ANALYSIS AND DESIGN}
\setcounter{section}{2}
\setcounter{subsection}{0}
\setcounter{figure}{0}
\setcounter{table}{0}

\subsection{Requirements Analysis}
The FootField platform is designed to serve three distinct user roles, each with specific functional requirements:

\begin{itemize}
    \item \textbf{Provider Admin (Super Admin):} The platform owner must be able to manage the entire SaaS ecosystem. This includes creating and managing subscription packages, overseeing tenant accounts, tracking global SaaS revenue, and responding to support tickets submitted by tenants.
    \item \textbf{Tenant Admin (Facility Owner):} Facility owners require tools to manage their day-to-day operations. Key requirements include managing football fields (pricing, types, maintenance status), viewing the booking schedule via an interactive time-grid to avoid conflicts, handling customer databases (CRM), calculating daily/monthly revenue, and using a mobile application to scan QR codes for rapid customer check-in.
    \item \textbf{Customer:} End-users need a seamless interface to browse available fields, check real-time availability, make bookings, and pay securely via VNPay. Additionally, they should be able to interact with an AI-powered chatbot to perform these tasks using natural language.
\end{itemize}

Non-functional requirements include high availability, responsive user interfaces across devices, and robust data isolation between tenants.

\subsection{System Architecture Design}

The system employs a Monolithic Hybrid Architecture. The core application runs as a single Node.js service containing both the backend API and the frontend assets. The database is hosted separately on a managed cloud service.

A critical architectural component is the AI Chatbot integration. The workflow is as follows:
\begin{enumerate}
    \item The customer inputs a request in natural language.
    \item The Node.js backend forwards the request to the Google Gemini API.
    \item Gemini analyzes the intent and, if necessary, triggers a "Function Call" (e.g., \texttt{check\_availability} or \texttt{create\_booking}).
    \item The backend executes the corresponding database queries and returns the results to the AI.
    \item The AI formulates a natural language response based on the query results and sends it back to the customer.
\end{enumerate}

\subsection{Database Design and Data Isolation}

The database is built on a relational model using MySQL. The core entities include \texttt{tenants}, \texttt{packages}, \texttt{users}, \texttt{fields}, \texttt{bookings}, \texttt{customers}, and \texttt{invoices}.

To implement the SaaS Multi-tenant architecture within a single database, the system utilizes a "Shared Database, Shared Schema" approach. This means all tenants store their data in the same tables. To ensure strict data isolation, almost every table (except global tables like \texttt{packages}) contains a \texttt{tenant\_id} foreign key. 

Every database query executed by the backend is scoped to the requesting user's \texttt{tenant\_id}. For example, when a Tenant Admin requests the list of bookings, the query is structured as \texttt{SELECT * FROM bookings WHERE tenant\_id = ?}. This approach is highly cost-effective and simplifies database maintenance while guaranteeing that tenants cannot access each other's data.